% titel: Blatt 7, A6: Ideale Gasgleichung - Druck, Volumen, Konzentration, Molvolumen
% Teil II: der durchgerechnete Weg.

\subsection*{Gegeben und gesucht}
a) $p=1.5$ atm, $V=300\ \mathrm{m^3}$, $T_1=10\,^\circ\mathrm{C}$,
Gesucht $p$ bei $30\,^\circ\mathrm{C}$. b) $V=650$ mL bei $160$ kPa,
Gesucht $V$ bei $200$ kPa (gleiche $T$). c) $c$ bei $0\,^\circ\mathrm{C}$
und $10^{-4}$ Pa. d) $V_m$ unter Standardbedingungen.

\subsection*{Ansatz}
Kenngr\"ossen der idealen Gasgleichung einsetzen und umstellen.

\subsection*{Rechnung}
\textbf{a)} Konstantes Volumen $\Rightarrow p\propto T$:
\[
  T_1=283.15\ \mathrm{K},\quad T_2=303.15\ \mathrm{K},
  \quad p_2 = p_1\,\frac{T_2}{T_1}
  = 1.5\cdot\frac{303.15}{283.15} = 1.61\ \mathrm{atm}.
\]

\textbf{b)} Konstante Temperatur (Boyle $p_1V_1=p_2V_2$):
\[
  V_2 = V_1\,\frac{p_1}{p_2} = 650\cdot\frac{160}{200}
  = 520\ \mathrm{mL}.
\]

\textbf{c)} $c = n/V = p/(RT)$, mit $p=10^{-4}$ Pa, $R=8.314$,
$T=273.15$ K:
\[
  c=\frac{10^{-4}}{8.314\cdot 273.15}
   = 4.40\times10^{-8}\ \mathrm{mol\,m^{-3}}
   = 4.40\times10^{-11}\ \mathrm{mol\,L^{-1}} .
\]

\textbf{d)} $V_m = RT/p$ unter Standardbedingungen ($p=1$ atm,
$T=273.15$ K):
\[
  V_m = \frac{0.0821\cdot 273.15}{1} = 22.4\ \mathrm{L\,mol^{-1}} .
\]

\begin{ergebnis}
a) $p=1.61$ atm. b) $V=520$ mL. c) $c=4.40\times10^{-11}\mathrm{mol\,L^{-1}}$
($=4.40\times10^{-8}\ \mathrm{mol\,m^{-3}}$). d) $V_m=22.4\ \mathrm{L\,mol^{-1}}$.
\end{ergebnis}

\begin{kontrolle}
Erw\"armung von $10$ auf $30\,^\circ\mathrm{C}$ (ca.\ 7\,\%) erh\"oht den
Druck proportional ($1.5 \to 1.61$ atm); h\"oherer Druck bei Boyle
verringert das Volumen ($650 \to 520$ mL). Das Standard-Molvolumen ergibt
den bekannten Wert $22.4$ L/mol.
\end{kontrolle}

\begin{merksatz}
Die ideale Gasgleichung umformen: konstant $V$ $\Rightarrow p\sim T$;
konstant $T$ $\Rightarrow pV=\text{const}$. Konzentration $=p/(RT)$,
Molvolumen $=RT/p$, Standardwert $22.4\ \mathrm{L\,mol^{-1}}$.
\end{merksatz}
